\chapter{Whitegape: A Destroy Board Authored From Nothing}
\label{ch:dtm}

\section{The Subject and Why It Is First}

Whitegape is a destroy-the-monument board for twenty players a side, authored
against no brief and taken from no composed arrangement. Every rectangle in its
plan was placed by its author, and the studio's only contribution to the
arrangement was to refuse the arrangements it would not accept.

That is what makes it the first of the two worked examples. Nothing about it is
inherited, so every decision in it is attributable, and the question this paper
asks at every stage has an answer that can be checked. Hushwater, in
\cref{ch:ctw}, begins from a board the studio produced, and the interesting
question there is the opposite one: what is left to decide once the arrangement
arrives already made.

\begin{figure}[htbp]
\centering
\includegraphics[width=\textwidth]{maps/whitegape-iso.png}
\caption{Whitegape as exported. The near half is team~A's: a stone quarry yard on
the rim, the monument on the floor of the pit cut into it, and a loading dock five
courses lower at the lip. \id{rot\_180} fans the whole of it to the far side.}
\label{fig:whitegape-iso}
\end{figure}

\section{The Board in One Sentence}

The board's author states it in the first line of the build specification: two
teams face each other across a chasm, each holding a stone-built quarry yard on
its own rim, with its monument on the floor of the pit cut into that yard, and
the only way over is a bridge the attackers build.

Everything downstream is a consequence of that sentence, and the sentence
is not the studio's. The studio has no opinion about what a map is of. It
will not suggest a quarry, it will not object to one, and it has no
representation in which \emph{limestone gorge} is distinguishable from
\emph{any other pair of rectangles at those coordinates}. What it has is a
refusal surface, and the sentence has to survive it.

The board is deliberately two kinds of place, and the seam between them is a
line a player can see. The back half, from $z=-100$ to $z=-64$, is grown
ground: open limestone fell, rolling, pinned in five places and left to the
solver everywhere else. The front half, from $z=-64$ to $z=-12$, is made
ground: one level stone yard with the pit sunk into it, a loading dock cut
five courses lower at the lip, a kiln on the dock and a works shed on the
yard. The retaining face between them carries two flights, and those two
flights are the only ways through it.

\section{The Plan: Four Rectangles and a Pillar}

The plan document is small. \Cref{tab:whitegape-plan} is the whole arrangement.

\begin{table}[htbp]
\caption{Whitegape's plan, in full. Rectangles are in cells; the grid is four
blocks to the cell.}
\label{tab:whitegape-plan}
\small
\begin{tabularx}{\textwidth}{L{2.6cm} L{2.4cm} L{3.4cm} Y}
\toprule
\textbf{Id} & \textbf{Role} & \textbf{Rect (cells)} & \textbf{What it is} \\
\midrule
\id{fell}   & \id{piece} & $[-7,-21,14,18]$ & the body of the board \\
\id{fell-w} & \id{piece} & $[-7,-25,2,4]$ & fell west of the spawn \\
\id{fell-e} & \id{piece} & $[-1,-25,8,4]$ & fell east of the spawn \\
\id{head}   & \id{spawn} & $[-5,-25,4,4]$ & the spawn shelf, one step higher \\
\id{chasm}  & zone & $[-7,-6,14,12]$ & the build region, over the whole gap \\
\id{destroyable-1} & objective & at $(-9,-50)$ & \id{pillar-3}, obsidian, float 4 \\
\bottomrule
\end{tabularx}
\end{table}

Two facts about that table carry most of the plan's weight.

The first is that the fell is tiled \emph{around} the spawn rather than drawn
under it. \id{PL4} refuses two pieces that overlap at different surfaces, and
the spawn shelf stands one block above the moor. The alternative was to
abandon the step, and the step is what \id{SP8} allows and what makes a spawn
read as a place rather than as a painted square, so the tiling is what gave
way. The board's back edge is therefore three rectangles where a reader looking
at the world sees one continuous fell.

The second is that the destroyable names no piece. Its \id{piece} field is
empty and its position is stated in blocks, because by the time it is placed
there is no rectangle under it: the pit is a shape sunk ten courses into the
yard, and the yard is a shape drawn over the fell. The plan level cannot
describe a hole, so the objective's home is described at the level that can,
and the plan states only where it is.

\subsection{Five Numbers, and Which of Them Are the Studio's}

The specification opens by naming five measurements and the rules they answer:
176 blocks from spawn to spawn along the lane, 42 blocks from a team's own
spawn to its goal, 56 blocks of width before the coast is cut, and a 24-block
void gap across the chasm. The rules cited are \id{GO1} (three to four goals
of separation), \id{GO3} (at least 85), and \id{GO4} (at least 40).

Every one of those bands is the studio's and every one of the numbers is the
author's. This is the division of labour in its plainest form: the studio does
not know what distance a board wants, and it knows exactly which distances it
will not accept. An author who does not read the bands first discovers them one
refusal at a time, which is the slowest way to learn a number that
\route{GET /api/rules} will print.

\section{The Coast: Editing an Outline the Studio Compiled}

Compiling the plan produces the board's ground, and the ground the plan
describes is rectangular. The compiled outline of the fell, which the studio
names \id{fell-24}, comes back as eight vertices, and three of its edges are
long straight runs down the flanks and across the lip.

A board that is left there reads as a rectangle from every camera. Whitegape's
author changes it in two ways, and the two are worth separating because they
are different instruments.

The first is deliberate: fourteen vertices are inserted and four are moved,
each one stating one place. A bay is pulled back into the west lip. The two
back corners are cut off. Nine points are run along the flanks and the lip, and
the lip's points stop short of the dock's own span, because a made edge is
straight and a grown one is not, and that difference is the whole argument of
the board.

The operations run in descending index. An insert shifts every index after it,
so a list written forwards and applied forwards addresses the wrong vertices
from the second operation onward. The specification replays each run in reverse
for it to come out forwards, which is a mechanical detail that costs a rebuild
to discover and nothing at all to state.

The second change is not deliberate in the same sense: a bend is applied to the
whole outline, with a wander of 2.5, a step of 11 and \id{side: out}. The
author chooses the parameters; the studio chooses every vertex the bend
produces. \id{side: out} is the one decision in it that is fully the author's,
and it is a safety decision rather than an aesthetic one: a bend that only
bloats cannot leave the yard or the dock hanging over void, and it cannot
narrow the strait out of \id{CT12}'s band of 15 to 40 blocks.

\section{Relief: Five Marks Over Half a Board}

\Cref{ch:relief} describes the height model. Whitegape uses it in the way that
model rewards and in one way it punished.

The rewarded use is that the relief is given half a board to spread over. Five
marks and three pushes are stated, and every one of them is in the back half or
along the lip. The front half is made ground, excluded from the solve entirely,
and takes no part in it. The consequence is that the solver is interpolating
between pins that are far apart, which is the condition under which it produces
ground rather than terraces.

The punished use was an earlier arrangement of the same marks. A single area
mark ran across the full width of the board at the foot of the retaining face,
which is the obvious way to state \emph{the fell meets the works at this
height}. It is also the way to flatten everything behind it, because a mark
pinned across the whole width leaves the solver nothing to interpolate
\emph{between}. The measurement that named it is the relief read: level 0.54
and largest field 0.21, both over the thresholds at which a surface is reported
to have no shape left. The fix was to replace the band with two aprons, one at
the foot of each flight, pinned only where a flight has to arrive at a definite
height.

That failure is worth stating in full because of its shape. Nothing refused it.
The board built, the export gate opened, and every picture of it looked like a
map. What said so was a number, and the number was available before the world
was built.

\begin{table}[htbp]
\caption{Whitegape's relief, stated in full. \id{base} 22, \id{reach} 0,
\id{step} 1, with a grain of amplitude 1.2 at scale 16.}
\label{tab:whitegape-relief}
\small
\begin{tabularx}{\textwidth}{L{3.0cm} L{1.8cm} Y}
\toprule
\textbf{Id} & \textbf{Kind} & \textbf{What it pins} \\
\midrule
\id{shelf}        & area & the spawn shelf at 25, one step over the moor \\
\id{gate-apron-w} & area & the foot of the west flight at 21 \\
\id{gate-apron-e} & area & the foot of the east flight at 21 \\
\id{liprim}       & area & the gorge lip at 19, either side of the dock \\
\id{cartway}      & line & the road off the fell down to the works gate \\
\midrule
\id{knott} & push & a rise on the fell \\
\id{slack} & push & a damp hollow the cart road skirts, $-5$ \\
\id{rigg}  & push & a low rise on the west flank, so the flank is not one gradient \\
\bottomrule
\end{tabularx}
\end{table}

\section{The Ground Finished by Its Angle}

The single most consequential decision about how Whitegape looks is that its
grown half is painted on the \id{slope} axis rather than on \id{height}.

A theme banded by height paints contour lines. Every column at $y=24$ takes the
same block whether it sits on a meadow or halfway up a cliff, and the result is
a board that reads flat from above however much relief is under it. A theme
banded by slope asks a different question of each column: not how high it is,
but how steeply the ground runs through it.

Whitegape's fell states three bands on that axis. Ground lying at up to
sixteen degrees takes grass over coarse dirt over dirt. From sixteen to
thirty-four degrees it takes a worn shoulder of coarse dirt and gravel over
dirt. Past thirty-four it takes bare rock, stone and andesite in a cell
pattern, two deep over stone. Each band is itself a layered material on the
\id{depth} axis, so the paint is a stack of stacks: slope chooses which
column profile applies, and depth walks it downward.

The cuts are not guessed. \route{GET .../incline?format=text} reports where the
board's ground actually lies, and the three bands are cut where that read said
the population divides. An author who picks band edges by eye produces a board
where one band is empty and another carries 90\% of the surface, and neither
the export nor any gate will mention it.

\section{The Made Half}

Fifteen shapes are added, and they are the front half of the board.

\begin{table}[htbp]
\caption{The fifteen added shapes, in the order the specification states them.}
\label{tab:whitegape-shapes}
\small
\begin{tabularx}{\textwidth}{L{2.4cm} L{1.9cm} L{1.9cm} Y}
\toprule
\textbf{Id} & \textbf{Type} & \textbf{Height} & \textbf{What it is} \\
\midrule
\id{yard} & polygon & level 24 & the terrace, 23 vertices, three notches cut into its front edge \\
\id{dock} & polygon & level 19 & the loading dock at the lip, front edge straight \\
\id{pit} & polygon & sink 10 & the quarry, sheer sides and a flat floor \\
\id{gate-w}, \id{gate-e} & polygon & flight & the two ways up through the retaining face \\
\id{stair-w}, \id{stair-e} & polygon & flight & yard down into the dock, filling two notches \\
\id{haul-ramp} & polygon & flight & the defenders' way down into the pit \\
\id{tramway} & polygon & flight & the attackers' way in: dock to pit floor, through the yard \\
\id{charge-ramp} & polygon & flight & yard up to the kiln's charging mouth \\
\id{revet-w}, \id{revet-m}, \id{revet-e} & polyline & level 27 & the retaining wall, in three runs \\
\id{kerb-w}, \id{kerb-e} & polyline & level 21 & the dock's parapet, split at the middle \\
\bottomrule
\end{tabularx}
\end{table}

Three things in that table are decisions rather than descriptions.

The pit is a \id{sink}, not a set of nested area marks. A sink cuts sheer and
leaves a flat floor. Nested marks at descending heights build a funnel, which is
what a quarry is not: a quarry has a floor and walls, and the floor is where
the monument stands.

The yard's three notches are sized to exactly the flights that fill them. A
flight drawn \emph{into} the yard would lose every one of its columns, because
the taller add wins the column outright and the yard stands five courses over
the dock. Cutting the notch out of the yard's own ring and dropping the flight
into it is the only arrangement in which both survive.

The revetment is in three runs rather than one. Its two gaps are exactly the
two gate flights, so the wall has ends, it has gates, and it has the fell on
one side and the works on the other. A single run with the flights drawn over
it would produce the same picture and a different fact: the wall would own
those columns and the flights would be decoration on top of a barrier.

\subsection{Made Ground Is Coursed, Not Bedded}

Whitegape's rock is stated once, as a six-band layered material on the depth
axis: two of stone, one of diorite, three of stone, one of andesite, four of
stone, two of diorite, repeating. Every cut face on the board reads those same
beds in that same order. The gorge walls, the quarry faces and the bank behind
the yard are one limestone rather than three separate decisions, and that is a
consequence of stating the material once and referring to it, not of any rule.

The made ground does not use it. The yard and the dock take a \id{wallRun}
instead: horizontal courses of stone brick five wide, stone three wide and
cobblestone two wide, striping along the perimeter. The distinction is
legible in the world at a glance, and it is the one place on the board where
\id{wallRun} is the right instrument. A run stripes along a face, which is what
masonry does and what a geological bed does not.

\section{The Kiln, Built Rather Than Emitted}

The board's one landmark is a lime kiln standing on the dock with its back
against the yard.

It is built out of rectangles on two made layers. That is a ruling rather than a
preference: the studio ships a library of parametric forms, and a tapered tower
off that library would have produced a kiln-shaped object in one call. The map
author's judgement, recorded during this run, is that a structure emitted from a
library does not read as a building, and that anything meant to be a building is
constructed by the author.

What the construction has to know is small and unforgiving. A layer holds one
span per column, so the base and the stack are two layers rather than two
heights on one; \id{SK9} declines the second span on a 200 and says so only in
\id{warnings}. The draw arch is drawn as the masonry \emph{around} the opening
rather than as a subtraction, because \id{SK13} reads a subtract as the board's
negative space and refuses any add that fills it. A rectangle is half-open, so
every piece needs a maximum strictly over its minimum or it draws no ground at
all, which \id{SK4} reports. The taper is five courses, each stepped in one
block, stated one rectangle at a time.

The kiln's material is brick over hardened clay, and it is the only brick on the
board. It is the one thing read from the far rim, and a landmark may not be the
ground it stands on.

\begin{figure}[htbp]
\centering
\begin{subfigure}{0.48\textwidth}
  \includegraphics[width=\linewidth]{maps/whitegape-made.png}
  \caption{Made ground. Every column a shape or a layer owns.}
\end{subfigure}\hfill
\begin{subfigure}{0.48\textwidth}
  \includegraphics[width=\linewidth]{maps/whitegape-structure.png}
  \caption{Structures. The spawn shell, the works shed, the kiln.}
\end{subfigure}
\caption{The made half, isolated by two of the studio's own read-backs. Neither
picture is drawn for this paper; both are routes on the running API.}
\label{fig:whitegape-made}
\end{figure}

\section{What the Board Says When It Is Read Back}

The export is not the end of the run. \Cref{ch:reading} describes the reads a
built world answers, and Whitegape's author takes roughly twenty-five of them on
every build. Three are worth reproducing, because each answers a question no
picture answers.

\subsection{The Heightmap Says Whether the Relief Solved}

The heightmap read prints one character per block, banded in single blocks above
a floor. Whitegape's back half runs from band 4 to band 9 across forty blocks
of the west flank and climbs to band \texttt{g} on the east, which is a
difference of twelve blocks over the width of the board. A flattened board
prints as a field of one character, and that is exactly what the earlier
full-width gate mark produced.

\subsection{The Slope Read Says Whether It Can Be Walked}

The slopes read prints the worst step from each block to a neighbour: a
full stop where the rise is nought or one, a colon where it is two, a hash where
it is three or more. Whitegape's fell is almost entirely full stops. The two
columns of hashes running down the middle of the back half are the spawn
shelf's own face, which is a wall by design, and the band of hashes across
$z=-85$ is the shelf's front. Everything an attacker crosses is walked.

\subsection{The Flow Read Says What the Board Does Not Use}

The flow read is the one that changes a board. It reports that 1952 of 9856
blocks, just under a fifth of the ground, sit off every route between the places
the board has, and it gives the coordinates: 848 blocks around $(-20, 68)$, 608
around $(16, -84)$, 432 around $(-24, -64)$.

The read's own wording is the point. Ground is dead because no journey passes
it, not because it is far out, and decorating it only means players look at it
on the way past. A fifth of Whitegape is ground nobody has a reason to stand on,
and the studio said so with coordinates rather than with an adjective.

\begin{figure}[htbp]
\centering
\begin{subfigure}{0.48\textwidth}
  \includegraphics[width=\linewidth]{maps/whitegape-traversability.png}
  \caption{Traversability. Where the worst step to a neighbour is a barrier.}
\end{subfigure}\hfill
\begin{subfigure}{0.48\textwidth}
  \includegraphics[width=\linewidth]{maps/whitegape-objectives.png}
  \caption{Objectives and their clearances.}
\end{subfigure}
\caption{Two read-backs of the finished world. Both are pictures of a
measurement rather than of the board.}
\label{fig:whitegape-reads}
\end{figure}

\section{What the Export Wrote}

The intent document states what the board is played for, and the export turns it
into a \id{map.xml} of 169 lines and a region folder of four files. The
gamemode is \id{dtm}. Two teams are declared at a maximum of twenty each, which
is the plan's \id{maxPlayers} carried through unchanged, and it is a cap per
team rather than for the board.

The monument is a three-block obsidian pillar floating four above the terrain.
The float is not a mistake and not a stylistic choice: a destroyable resting on
the ground is trivially covered, and floating it is what PGM has always done.
This is the one claim in this chapter that the repository cannot settle. Neither
the corpus nor the studio's own code says what a destroyable needs around it;
the map author does, and the answer is recorded rather than derived.

\section{The Division of Labour on This Board}

\Cref{tab:whitegape-labour} is the chapter in one table.

\begin{table}[htbp]
\caption{Whitegape, stage by stage: what the studio settled and what was left
to its author.}
\label{tab:whitegape-labour}
\small
\begin{tabularx}{\textwidth}{L{2.6cm} Y Y}
\toprule
\textbf{Stage} & \textbf{\studio{The studio}} & \textbf{\agent{The author}} \\
\midrule
Arrangement & \id{PL4}: no two pieces overlapping at different surfaces.
  \id{GO1}, \id{GO3}, \id{GO4}: the distance bands &
  four rectangles, one zone, one objective, and the five numbers inside
  the bands \\
Coast & every vertex the bend produced; the compiled eight-vertex outline &
  fourteen inserts, four moves, and \id{side: out} \\
Relief & the solved field between the pins; the level and field measures
  that named the flattening &
  five marks, three pushes, and the decision to exclude the front half \\
Paint & which block lands on which cell; the incline report &
  four themes, the slope axis, and where the three bands cut \\
Shapes & the taller add wins; \id{SK4}, \id{SK9}, \id{SK13} &
  fifteen outlines, and a notch sized to each flight \\
Structures & storey heights, roof pitch, which side the door takes &
  the kiln as rectangles, and the refusal to emit one \\
Dressing & the legal seats and every decline &
  ten styles and nineteen props \\
Objectives & the region, the filter, the monument's own blocks &
  what the board is played for, and the float \\
\bottomrule
\end{tabularx}
\end{table}
